\chapter{Leitungen und Kabel}
\label{chap:05-leitungen-und-kabel}

\section{Einordnung und Modellierungsziele}

Leitungen und Kabel verbinden Betriebsmittel und bestimmen durch ihre
Serienimpedanzen und Queradmittanzen den Spannungsfall, die Verluste,
die Kurzschlussstromhöhe und das Erdschlussverhalten.

Freileitungen und Kabel unterscheiden sich unter anderem in ihrer Geometrie,
ihrem Isolationssystem, ihrer Wärmeabfuhr und ihrer kapazitiven Kopplung
zur Erde. Die Modellierung muss deshalb zum betrachteten Phänomen passen.

Für eine erste Netzberechnung wird häufig ein einphasiges Ersatzmodell
je Phase verwendet. Bei unsymmetrischen Fehlern und Erdschlussberechnungen
müssen zusätzlich die Sequenzimpedanzen und die Rückstrompfade berücksichtigt werden.

\section{Längs- und Querparameter}

Die längenbezogenen Serienparameter einer Leitung werden durch
\[
  Z'=R'+j\omega L'
\]
beschrieben. Die Queradmittanz lautet
\[
  Y'=G'+j\omega C'.
\]

Dabei bezeichnen:
\begin{itemize}
  \item \(R'\): längenbezogener Wirkwiderstand;
  \item \(L'\): längenbezogene Induktivität;
  \item \(G'\): längenbezogener Ableitwert;
  \item \(C'\): längenbezogene Kapazität;
  \item \(\omega=2\pi f\): Kreisfrequenz.
\end{itemize}

Für eine Leitung der Länge \(l\) ergibt das konzentrierte Modell
\[
  Z=Z'l,\qquad Y=Y'l.
\]

Diese Gleichungen setzen voraus, dass die Parameter entlang der Leitung
hinreichend gleichmäßig sind. Bei langen Leitungen oder wenn die Queradmittanz
wesentlich ist, ist eine verteilte Leitungsbeschreibung genauer.

\section{Ohmscher Widerstand}

\subsection{Grundgleichung}

Für einen homogenen Leiter mit Länge \(l\), Querschnitt \(A\) und
spezifischem Widerstand \(\rho\) gilt
\[
  R=\rho\frac{l}{A}.
\]

Der Widerstand wächst proportional mit der Länge und sinkt bei größerem
Querschnitt. Die Materialeigenschaften gehen über \(\rho\) ein.

Bei Wechselstrom können Skin- und Proximity-Effekte den wirksamen Widerstand
gegenüber dem Gleichstromwert erhöhen. Bei mehradrigen Kabeln und dicht
benachbarten Leitern ist deshalb nicht immer der reine Gleichstromwiderstand
für Verlust- und Spannungsfallberechnungen ausreichend.

\subsection{Temperaturabhängigkeit}

In einem begrenzten Temperaturbereich lässt sich der Widerstand näherungsweise
mit
\[
  R_\vartheta=R_{\vartheta_0}
  \left[1+\alpha(\vartheta-\vartheta_0)\right]
\]
korrigieren. \(\alpha\) ist der materialabhängige Temperaturkoeffizient.

Die tatsächliche Leitertemperatur ist nicht nur eine Eingangsgröße:
Sie stellt sich aus Verlustleistung, thermischem Widerstand der Umgebung,
Verlegebedingungen und Wärmeabfuhr ein. Für die Bemessung ist deshalb eine
thermische Betrachtung erforderlich.

\section{Induktivität und Reaktanz}

Die Induktivität hängt von Leitergeometrie, Leiterabständen, Rückstrompfaden
und magnetischen Eigenschaften der Umgebung ab. Für die Serienreaktanz gilt
\[
  X_L=\omega L.
\]

Die Reaktanz wächst bei gleichbleibender Induktivität linear mit der Frequenz.
Bei Wechselstromnetzen beeinflusst sie den Spannungsfall und den Blindleistungsfluss.

Bei symmetrischen Drehstromleitungen wird die gegenseitige magnetische Kopplung
der Phasen in den Ersatzparametern berücksichtigt. Bei unsymmetrischer Geometrie
oder asymmetrischer Leitungsführung kann eine einfache Einphasenbetrachtung
ungenau sein.

\section{Kapazität und Ableitwert}

\subsection{Kapazitive Kopplung}

Zwischen Leitern sowie zwischen Leitern und geerdeten Teilen entstehen
elektrische Felder. Daraus folgt eine Kapazität. Der kapazitive Strom eines
einfachen Zweiges beträgt
\[
  I_C=\omega C U.
\]

Die verwendete Spannung \(U\) muss zur Definition der Kapazität passen.
Eine Kapazität zwischen zwei Phasen wird nicht ohne Umrechnung wie eine
Kapazität einer Phase gegen Erde behandelt.

Der Ableitwert \(G\) beschreibt reale dielektrische Verluste und Leckpfade.
Bei vielen vereinfachten Netzberechnungen ist er gegenüber dem kapazitiven
Anteil klein, kann bei speziellen Isolations- und Betriebsbedingungen aber
relevant sein.

\subsection{Kabel gegenüber Freileitungen}

Kabel besitzen aufgrund der geringen Leiterabstände und der hohen relativen
Permittivität ihrer Isolierung typischerweise eine deutlich größere Kapazität
je Längeneinheit als Freileitungen.

Damit steigt bei gleicher Spannung und Frequenz der kapazitive Ladestrom
mit der Kabellänge. In Mittelspannungsnetzen kann dieser Zusammenhang das
Erdschlussverhalten wesentlich beeinflussen.

\section{Leitungsmodelle}

\subsection{Kurze Leitung}

Bei einer kurzen Leitung kann die Queradmittanz für bestimmte Fragestellungen
vernachlässigt werden. Das Ersatzmodell besteht dann aus einer Serienimpedanz
\[
  Z=R+jX.
\]

Dieses Modell ist besonders für überschlägige Spannungsfall- und
Kurzschlussberechnungen nützlich. Es bildet jedoch keine verteilten
kapazitiven Ströme ab.

\subsection{Nominales Pi-Modell}

Bei einer Leitung mittlerer Länge wird die Queradmittanz häufig auf beide
Leitungsenden aufgeteilt. Das nominale Pi-Modell enthält die Serienimpedanz
\(Z\) und zwei Querzweige mit jeweils \(Y/2\).

Die Näherung bildet den Einfluss der Leitungskapazität besser ab als das
reine Serienmodell. Sie ist jedoch nicht identisch mit einer vollständigen
verteilten Leitungsrechnung.

\subsection{Verteilte Parameter}

Bei einer Leitung mit gleichmäßig verteilten Parametern werden Spannung
und Strom entlang der Leitung durch die Leitungsgleichungen beschrieben.
Mit den längenbezogenen Größen \(Z'\) und \(Y'\) ist die Ausbreitungskonstante
\[
  \gamma=\sqrt{Z'Y'}
\]
und der Wellenwiderstand
\[
  Z_W=\sqrt{\frac{Z'}{Y'}}.
\]

Die Zweitorbeziehung für eine homogene Leitung der Länge \(l\) kann in der Form
\[
  \begin{bmatrix}
    U_1\\ I_1
  \end{bmatrix}
  =
  \begin{bmatrix}
    \cosh(\gamma l) & Z_W\sinh(\gamma l)\\
    \frac{1}{Z_W}\sinh(\gamma l) & \cosh(\gamma l)
  \end{bmatrix}
  \begin{bmatrix}
    U_2\\ I_2
  \end{bmatrix}
\]
geschrieben werden, sofern die Stromrichtungen und die Zuordnung der
Endgrößen konsistent definiert sind.

Dieses Modell ist insbesondere für lange Leitungen relevant. Die genaue
Anwendung erfordert eine einheitliche Vorzeichenkonvention für die Ströme.

\section{Spannungsfall bei Drehstrom}

Für eine kurze, symmetrisch belastete Drehstromleitung kann der Spannungsfall
näherungsweise durch
\[
  \Delta U\approx\sqrt{3}\,I\,l
  \left(R'\cos\varphi+X'\sin\varphi\right)
\]
abgeschätzt werden.

Hier sind \(R'\) und \(X'\) die längenbezogenen Serienparameter je Phase,
\(I\) der Leiterstrom und \(l\) die einfache Leitungslänge. Die Formel
verwendet die übliche Näherung für kleine Spannungsänderungen und einen
definierten Leistungsfaktorwinkel \(\varphi\).

Bei induktiver Last ist \(\sin\varphi\) positiv. Bei führendem Leistungsfaktor
kann der Blindleistungsanteil das Vorzeichen wechseln. Eine kapazitive
Einspeisung kann daher den Spannungsfall verringern oder eine Spannungsanhebung
verursachen.

Für größere Spannungsänderungen, lange Leitungen oder bedeutende Queradmittanzen
ist eine Phasorrechnung oder ein vollständiges Leitungsmodell zu verwenden.

\section{Leitungsverluste}

Die ohmschen Verluste eines symmetrischen Drehstromsystems sind näherungsweise
\[
  P_{\mathrm V}=3I^2R.
\]

Dabei ist \(R\) der Widerstand eines Leiters über die betrachtete Leitungslänge.
Zusätzliche Verluste können in metallischen Kabelschirmen, Bewehrungen,
Erdungsleitern und anderen leitfähigen Komponenten auftreten.

Für eine gegebene übertragene Wirkleistung \(P\) und einen gegebenen
Leistungsfaktor gilt
\[
  I=\frac{P}{\sqrt{3}\,U\,\cos\varphi}.
\]
Ein höheres Spannungsniveau reduziert daher bei gleicher Leistung und gleichem
Leistungsfaktor den Strom und damit die ohmschen Verluste, sofern die übrigen
Parameter vergleichbar bleiben.

\section{Thermische Strombelastbarkeit}

Die zulässige Dauerstrombelastbarkeit ist eine thermische Bemessungsgröße.
Sie hängt insbesondere ab von:
\begin{itemize}
  \item Leiterwerkstoff und Leiterquerschnitt;
  \item zulässiger Leitertemperatur und Isolationssystem;
  \item Umgebungstemperatur und Verlegeart;
  \item Häufung mehrerer belasteter Kabel;
  \item thermischem Widerstand des Bodens bei Erdkabeln;
  \item Verlusten in Schirmen und metallischen Umhüllungen;
  \item Lastprofil und thermischen Randbedingungen.
\end{itemize}

Eine Querschnittswahl allein anhand des ohmschen Widerstands ist deshalb
nicht ausreichend. Zusätzlich sind Spannungsfall, Kurzschlussfestigkeit,
Schutzbedingungen, mechanische Anforderungen und die einschlägigen
Bemessungsregeln zu berücksichtigen.

\section{Kurzschlussbeanspruchung von Leitern}

Während eines Kurzschlusses kann der Strom wesentlich höher als der
Dauerstrom sein. Für eine vereinfachte adiabatische Betrachtung wird häufig
eine Beziehung der Form
\[
  S=\frac{I_k\sqrt{t}}{k}
\]
verwendet, wobei \(S\) den Leiterquerschnitt, \(I_k\) den wirksamen
Kurzschlussstrom, \(t\) die Fehlerdauer und \(k\) einen vom Werkstoff,
Isolationssystem und den zulässigen Anfangs- und Endtemperaturen abhängigen
Faktor bezeichnet.

Diese Beziehung darf nur mit einem zum Leiter und zur Isolierung passenden
Faktor angewendet werden. Bei stark zeitabhängigen Strömen, mehreren
Fehlerbeiträgen oder nichtadiabatischen Bedingungen sind die jeweils
zutreffenden Verfahren und Normen heranzuziehen.

\section{Kabelschirme und Erdung}

Metallische Schirme und Bewehrungen erfüllen elektrische, mechanische und
feldsteuernde Aufgaben. Ihre Erdung und Verbindung beeinflussen:
\begin{itemize}
  \item die Sicherheit bei Isolationsfehlern;
  \item die Höhe und Verteilung von Schirmströmen;
  \item die induzierten Spannungen;
  \item die Nullsystemimpedanz;
  \item die thermische Beanspruchung bei Erd- und Kurzschlüssen.
\end{itemize}

Beidseitige Schirmerdung kann unter bestimmten Bedingungen geschlossene
Strompfade für induzierte Ströme erzeugen. Einseitige Erdung reduziert solche
Schirmströme möglicherweise, kann aber am nicht geerdeten Ende höhere
induzierte Spannungen verursachen. Die Wahl der Schirmerdung ist deshalb
eine Systementscheidung und keine allgemeine Regel, die unabhängig von
Kabellänge, Spannungsebene und Fehlerfall getroffen werden kann.

\section{Übungsaufgaben}

\begin{enumerate}
  \item Berechnen Sie den Gleichstromwiderstand eines Kupferleiters mit
        vorgegebenem Querschnitt und einer Länge von \(250\,\mathrm m\).
        Berücksichtigen Sie anschließend eine Temperaturerhöhung.
  \item Ein Kabel besitzt eine Kapazität von \(0{,}25\,\mathrm{\mu F/km}\)
        je Phase gegen Erde. Berechnen Sie den kapazitiven Strom für eine
        Leitungslänge von \(8\,\mathrm{km}\) bei \(20\,\mathrm{kV}\) und
        \(50\,\mathrm{Hz}\). Geben Sie an, welche Spannung Sie einsetzen.
  \item Erläutern Sie die Unterschiede zwischen Serienmodell, nominalem
        Pi-Modell und verteiltem Leitungsmodell.
  \item Warum kann die Schirmerdung die Erdschlussberechnung beeinflussen?
  \item Welche Informationen fehlen, wenn ein Kabel allein anhand des
        Querschnitts thermisch bemessen werden soll?
\end{enumerate}